% =====================================================================
% CAPITULO 2 - ESTADO DEL ARTE
% =====================================================================
% Criterio: de cada familia se dice que aporta, que cuesta y SI ESTE TRABAJO
% LA HA PROBADO O NO. La seccion de defectos de los datasets pesa tanto como
% la taxonomia, porque es donde enlaza con los hallazgos propios.
%
% ESTILO (31-ago): sin negritas en el texto corrido, sin incisos entre rayas
% salvo cuando la frase los pida, y sin florituras.
% =====================================================================

Este capítulo sitúa el trabajo en su campo con una restricción declarada: de cada
familia de modelos se dice qué aporta, qué cuesta y si este trabajo la ha
contrastado o no. De las siete familias que se repasan, tres se han contrastado
experimentalmente y cuatro no. Enumerarlas todas sin distinguirlo produciría un
estado del arte más lucido y menos honesto.

La segunda mitad del capítulo trata de los conjuntos de datos públicos, de cómo
se evalúan los modelos sobre ellos y de sus defectos conocidos. Ocupa tanto
espacio como la taxonomía de modelos, y esa proporción es deliberada: dos de los
hallazgos centrales de esta memoria son mediciones de esos defectos, no mejoras
de un modelo.

\section{Detección de intrusiones: de las firmas al aprendizaje}
\label{sec:sota-contexto}

Un sistema de detección de intrusiones en red analiza el tráfico en busca de
actividad maliciosa. Durante décadas lo hizo mediante firmas: reglas que
describen patrones conocidos, como una cadena concreta en la carga útil o una
secuencia de puertos, que un motor contrasta contra cada paquete.

El enfoque tiene una virtud considerable, que es la precisión, y un límite
estructural: solo detecta lo que ya se ha visto y catalogado. Frente a variantes
polimórficas, tráfico cifrado o ataques nuevos, una base de firmas queda ciega
por construcción, y actualizarla exige que alguien haya analizado antes la
amenaza \cite{Swarnkar2023}.

De ahí el giro hacia el aprendizaje automático, que sustituye la regla escrita a
mano por un modelo inducido de los datos y aspira a generalizar a variantes no
vistas. El giro trae sus propios problemas, entre ellos los falsos positivos, la
opacidad del modelo y la dependencia de datos etiquetados, y esta memoria añade
uno más, que es su tesis: la dificultad de medir si el modelo ha aprendido lo
que se cree.

Hay además dos formas de plantear la tarea, porque la literatura
las mezcla y los resultados no son comparables. La detección basada en firmas o
en clasificación supervisada parte de ataques etiquetados y aprende a
reconocerlos; la detección de anomalías se entrena únicamente con tráfico
benigno y marca lo que se desvía. La primera obtiene resultados mucho mejores
sobre los conjuntos públicos, pero por construcción solo reconoce lo que se
parece a lo que vio. La segunda es la que aspira a detectar ataques de día cero,
a costa de una tasa de falsos positivos mucho más alta. Este trabajo se centra en
la primera, y emplea la segunda solo como contraste.

Ese contraste está medido y vale la pena anticiparlo, porque acota cuánto valen las
cifras del resto de la memoria. La suposición de que la mayor parte del tráfico
es benigno permite entrenar sin etiquetas de ataque, y la
sección~\ref{sec:anomalia} lo hace con dos métodos clásicos, un Isolation Forest
y un autoencoder, ajustados únicamente sobre tráfico benigno y evaluados con el
área bajo la curva ROC para no depender de un umbral. El resultado es que
retirar las etiquetas de ataque cuesta entre veinte y cincuenta puntos según el
régimen. Esa es la distancia entre lo que se mide sobre un conjunto etiquetado y
lo que un detector encontraría ante una amenaza nueva.

\section{Taxonomía de enfoques}
\label{sec:sota-taxonomia}

Las familias que siguen se ordenan de menor a mayor complejidad estructural. La
tabla~\ref{tab:taxonomia} las resume y las subsecciones desarrollan cada una.

\begin{table}[htbp]
  \centering
  \caption[Familias de modelos y cuáles se han contrastado]{Familias de modelos
  para detección de intrusiones y ejemplos representativos de cada una.}
  \label{tab:taxonomia}
  \begin{tabular}{llc}
    \toprule
    \textbf{Familia} & \textbf{Ejemplos} & \textbf{Contrastada} \\
    \midrule
    Ensembles sobre flujo   & Random Forest, XGBoost   & Sí \\
    Tabular profundo        & TabNet                   & No \\
    Grafos                  & E-GraphSAGE              & No \\
    Payload como imagen     & CNN 2D sobre bytes       & No \\
    \emph{Transformers}     & ET-BERT, DeBERTa         & Sí \\
    Espacios de estados     & Mamba-2                  & No \\
    Híbridos multimodales   & Fusión tardía, atención cruzada & Sí \\
    \bottomrule
  \end{tabular}
\end{table}

De las siete familias, este trabajo ha contrastado experimentalmente tres. Las
otras cuatro se describen porque delimitan el campo, y en cada caso se indica por
qué quedaron fuera. Cada subsección expone qué aporta la familia y a qué coste.

\subsection{Ensembles sobre características de flujo}

Random Forest, XGBoost \cite{Chen2016} y LightGBM operan sobre vectores de
características agregadas del flujo, como la duración, el número de paquetes o
los estadísticos de tamaño, y constituyen el estándar de facto del campo. Su
rendimiento sobre los conjuntos de datos públicos es difícil de superar y su
coste de inferencia es mínimo, lo que los hace además desplegables en línea.

El motivo de su éxito es que las características de flujo son de dimensión baja,
están disponibles siempre y son razonablemente informativas: un ataque de
denegación de servicio produce conexiones cortas y numerosas, y una fuerza bruta
produce conexiones de tamaño casi idéntico. Los ensembles de árboles capturan
bien ese tipo de estructura, con umbrales sobre variables individuales, y toleran
sin dificultad la mezcla de escalas.

Su límite es la miopía contextual. Una característica agregada no contiene el
contenido, así que un ataque cuya firma esté en los bytes queda fuera de su
alcance. Tampoco captura relaciones entre flujos, porque cada conexión se
clasifica de forma independiente.

Esta familia es la referencia con la que se contrasta en el
capítulo~\ref{ch:generalizacion}, y el resultado se puede adelantar, porque
matiza la afirmación de supremacía. Dentro del día los ensembles ganan a un
perceptrón multicapa sobre la vista tabular, pero por márgenes que van de
milésimas a $4.7$ puntos, y esa ventaja no sobrevive al cambio de condiciones:
al transferir entre días, los tres modelos colapsan por igual salvo cuando se les
da acceso a las direcciones IP, que no es una alternativa aceptable sino la
medición de un error: la dirección identifica al atacante de este laboratorio y
no describe el ataque, y lo que mejora con ella no es la detección sino
la memorización. La sección~\ref{sec:fuga} lo cuantifica.

\subsection{Modelos tabulares profundos}

TabNet \cite{Arik2021} aplica atención secuencial a datos tabulares, buscando la
potencia de una red profunda con la interpretabilidad de una selección explícita
de características. En cada paso de decisión el modelo elige qué subconjunto de
variables mirar, lo que permite explicar sus predicciones sin recurrir a métodos
externos. Aplicado a la detección de intrusiones, TabNet-IDS
\cite{Rodriguez2023} reporta una exactitud media del 97.03\,\% en CIC-IDS2017 y
del 95.58\,\% en CSE-CIC-IDS2018.

No se ha probado en este trabajo. Sobre los datos aquí empleados la comparación
tendría poco valor informativo, porque el capítulo~\ref{ch:regimenes} muestra que
la vista tabular ya satura y un modelo tabular más potente no tendría margen que
recoger. La interpretabilidad que aporta TabNet sí sería útil, pero el trabajo la
obtiene por otra vía, mediante saliencia sobre el modelo de bytes.

\subsection{Redes neuronales de grafos}

Las redes sobre grafos, entre ellas E-GraphSAGE \cite{Lo2022}, modelan el tráfico
como un grafo de comunicaciones entre \emph{hosts}: los nodos son máquinas, las
aristas son flujos, y el modelo aprende representaciones que incorporan el
vecindario de cada nodo. Detectan así anomalías topológicas, es decir, patrones
en quién habla con quién que ningún flujo aislado revela.

Es la familia natural para el movimiento lateral, que por definición no vive en
un flujo sino en la estructura de las comunicaciones. Su coste es la gestión de
estado: construir y mantener el grafo en línea es sustancialmente más complejo
que clasificar flujos independientes, y la evaluación exige decidir cómo se
segmenta el grafo en ventanas.

No se ha probado, y la ironía salta a la vista, porque es precisamente el enfoque
que capturaría el caso que la sección~\ref{sec:limite} identifica como el límite
de este trabajo. El reconocimiento de la Infiltration es un abanico de conexiones
sin carga útil desde un único origen hacia 612 destinos, es decir, un patrón
puramente topológico. Queda como la línea de trabajo futuro con mejor retorno
esperado, y se propone en la sección~\ref{sec:futuro} por ese motivo y no por
completitud.

\subsection{Payload como imagen: redes convolucionales}

Una línea extendida, que arranca con la clasificación de tráfico de \emph{malware}
mediante representación aprendida \cite{Wang2017}, convierte los bytes de la carga
útil en una imagen bidimensional, rasterizándolos en una matriz, y aplica una red
convolucional 2D. La motivación es práctica: permite reutilizar arquitecturas
maduras de visión por computador, con sus pesos preentrenados y sus técnicas de
aumento de datos.

La objeción es de fondo y condicionó el diseño de este trabajo. La rasterización
destruye la causalidad temporal. Un flujo de red es una secuencia en la que el
byte $n$ precede al $n+1$; al plegarlo en una matriz, bytes distantes en el
tiempo quedan adyacentes en el espacio y la convolución 2D los trata como
vecinos. La estructura que se le pide aprender al modelo no es la que tienen los
datos.

Por eso el byte-CNN de esta memoria opera sobre la secuencia unidimensional, con
convoluciones de varios tamaños en paralelo. La decisión resultó acertada por una
razón que solo se supo después: la sección~\ref{sec:saliency} muestra que la
firma se concentra en los primeros 48 bytes, y una convolución con agrupación
máxima localiza esa región, mientras que una rasterización la habría repartido
entre filas no contiguas.

\subsection{Transformers y modelos semánticos}

ET-BERT \cite{ETBERT} y las variantes de DeBERTa aplican preentrenamiento y
autoatención al tráfico, tratando el paquete como una secuencia de \emph{tokens}
con semántica propia. El preentrenamiento se hace sobre grandes volúmenes de
tráfico sin etiquetar, con objetivos autosupervisados análogos a los del
procesamiento de lenguaje natural, y el modelo resultante se ajusta después a la
tarea concreta con muchos menos ejemplos etiquetados.

Es la familia con mejores resultados publicados en clasificación de tráfico
cifrado y malware ofuscado. Su coste principal es la autoatención, cuadrática en
la longitud de la secuencia, lo que obliga a truncar los flujos o a segmentarlos.
A ello se añade el coste del preentrenamiento, que no siempre se contabiliza al
comparar.

Se ha probado, con una restricción importante, que se declara aquí: se contrasta
la arquitectura, es decir, un codificador con autoatención a presupuesto de
parámetros comparable, y no el preentrenamiento, que es buena parte de lo que
aporta ET-BERT. El resultado, en la sección~\ref{sec:sota-contraste}, es que el
\emph{transformer} nunca gana y cuesta entre 36 y 49 veces más tiempo. Esa
conclusión vale para esta comparación y no refuta el enfoque en general.

\subsection{Modelos de espacio de estados}

Las arquitecturas basadas en espacios de estados, como Mamba-2, prometen lo
mismo que un \emph{transformer} con coste lineal en la longitud de la secuencia
en lugar de cuadrático. Aplicadas al tráfico permiten operar directamente sobre
los bytes sin \emph{tokenización} ni preentrenamiento, lo que elimina dos
decisiones de diseño que en los modelos semánticos condicionan mucho el resultado
\cite{MambaNetBurst2026}.

No se ha probado. Es la familia más reciente de las siete, su ecosistema de
implementaciones es todavía joven, y queda fuera del alcance de este trabajo.

\subsection{Arquitecturas híbridas y fusión multimodal}

Es la familia en la que se sitúa esta memoria. Combina una rama que procesa el
contenido con otra que procesa las características de flujo, y la cuestión de
diseño es dónde y cómo fusionarlas. Caben al menos tres respuestas: concatenar
las entradas antes de procesarlas, fundir las representaciones intermedias que
cada rama produce, o dejar que una rama atienda a la otra mediante atención
cruzada \cite{TwoBranch2026}. Otras propuestas recientes combinan un autoencoder
con un clasificador clásico \cite{HybridAutoencoder2024}.

Este trabajo emplea fusión tardía, en la que cada rama construye su
representación por separado y se concatenan antes de la capa de decisión, y
contrasta esa elección con la atención cruzada del estado del arte. El resultado
es un empate con 2.3 veces menos parámetros
(sección~\ref{sec:sota-contraste}).

Falta un matiz que la literatura no siempre explicita. La fusión
ingenua, concatenando las entradas crudas, es peor que cualquiera de las dos
ramas por separado, porque las 256 dimensiones del histograma diluyen los cuatro
escalares de flujo y el modelo acaba ignorándolos. La diferencia entre fusionar
entradas y fusionar representaciones no es un detalle de implementación.

\section{Cómo se evalúa en la literatura}
\label{sec:sota-evaluacion}

Los modelos anteriores se comparan casi siempre con el mismo protocolo, y ese
protocolo es el objeto de buena parte de esta memoria. Conviene por tanto
describirlo antes de pasar a los conjuntos de datos.

\subsection*{La partición}

Lo habitual es tomar un conjunto de datos etiquetado, mezclar todos sus flujos y
partirlos aleatoriamente en entrenamiento y prueba, o bien aplicar validación
cruzada de $k$ particiones. El reparto es aleatorio a nivel de flujo, sin
considerar de qué host procede cada uno ni en qué momento se capturó.

Esa elección tiene una consecuencia que rara vez se discute: flujos de la misma
conexión lógica, del mismo host y del mismo minuto acaban repartidos entre
entrenamiento y prueba. El modelo no tiene que generalizar a una situación nueva,
solo a otros ejemplos de la misma situación. El capítulo~\ref{ch:regimenes} mide
qué queda de las diferencias entre representaciones bajo ese protocolo, y la
respuesta es que no queda ninguna.

\subsection*{Las métricas}

Se reportan mayoritariamente exactitud, precisión, \emph{recall} y F1, y con
frecuencia solo la primera. Sobre conjuntos donde la clase benigna supera el
95\,\% de los ejemplos, la exactitud es poco informativa, porque un clasificador
que responda siempre «benigno» la maximiza. El problema es conocido, y aun así
las cifras de exactitud siguen apareciendo como resultado principal.

Menos frecuente es el uso de métricas independientes del umbral, como el área
bajo la curva ROC o la curva de precisión-\emph{recall}. La
sección~\ref{sec:cross-dia} muestra un caso en que esa ausencia cambia por
completo la conclusión: un modelo que con umbral fijo parece no transferir en
absoluto conserva un AUC superior a $0.97$: lo que falla es el
umbral y no la representación.

Prácticamente ausente está la calibración, es decir, la pregunta de si las
probabilidades que el modelo emite se corresponden con las frecuencias reales.
Sobre un único conjunto de datos la calibración importa poco, porque el umbral se
ajusta una vez; entre dominios distintos es justamente lo primero que se rompe.

\subsection*{Lo que casi nunca se mide}

Tres comprobaciones aparecen rara vez en la literatura revisada y son las que
este trabajo emplea como eje del capítulo~\ref{ch:generalizacion}.

La primera es la evasión activa, es decir, comprobar qué queda del detector
cuando el atacante modifica deliberadamente la parte de la entrada en la que el
modelo se apoya. Sin esa comprobación no se puede distinguir una firma del ataque
de una firma de la herramienta con que se ejecutó.

La segunda es la interpretabilidad aplicada al diagnóstico, y no solo a la
presentación. Saber qué parte de la entrada usa el modelo permite anticipar qué
lo romperá, y en este trabajo fue el paso que motivó todos los demás.

La tercera es la transferencia entre dominios, entendida como entrenar en unas
condiciones y evaluar en otras distintas. Cuando aparece, suele hacerse entre
particiones del mismo conjunto de datos, lo que como muestra la
sección~\ref{sec:fuga} no basta para detectar los atajos que el propio conjunto
introduce.

\section{Conjuntos de datos públicos y sus defectos}
\label{sec:sota-datasets}

Esta sección es la que más conecta con el resto de la memoria, porque los
defectos que describe no son una advertencia general: son el objeto de dos de los
hallazgos de este trabajo.

\subsection*{Los conjuntos disponibles}

CIC-IDS2017 y CSE-CIC-IDS2018 \cite{Sharafaldin2018, CSECICIDS2018} son los
conjuntos de referencia del campo. Se generan en un laboratorio controlado, con
perfiles de tráfico benigno sintetizados a partir del comportamiento de usuarios
reales y ataques ejecutados con herramientas conocidas y públicas, como Hulk,
GoldenEye, Slowloris, LOIC o Patator. Se distribuyen tanto en capturas de red en
bruto como en ficheros de características de flujo ya etiquetadas, extraídas con
la herramienta CICFlowMeter.

Esa doble distribución explica buena parte de la literatura. Trabajar sobre los
ficheros de características es mucho más barato, porque evita procesar centenares
de gigabytes de capturas, y por eso la mayoría de los trabajos publicados los
usan. El precio es que se pierde el acceso al contenido y se hereda cualquier
defecto que la herramienta de extracción introduzca. Este trabajo parte de las
capturas en bruto precisamente para evitar las dos cosas.

UNSW-NB15 \cite{Moustafa2015} sigue una filosofía parecida con otra taxonomía de
ataques y un generador de tráfico distinto.
Malware-Traffic-Analysis.net \cite{MalwareTrafficAnalysis} aporta capturas de
infecciones reales, sin sintetizar, a cambio de no venir etiquetadas de forma
sistemática. Ninguno de los dos se ha usado aquí; la
sección~\ref{sec:futuro} explica por qué el segundo sería el complemento más
valioso.

\subsection*{Los cuatro defectos documentados}

Engelen, Rimmer y Joosen \cite{Engelen2021} auditaron CIC-IDS2017 contra sus
capturas originales y documentaron cuatro problemas. Los cuatro se reproducen en
CSE-CIC-IDS2018, que comparte metodología de generación.

\begin{enumerate}
  \item Errores de etiquetado. Flujos marcados con la clase equivocada por
        desajustes entre las ventanas declaradas y la actividad real.
  \item Defectos de la herramienta de extracción de características, que produce
        valores incorrectos en algunas columnas.
  \item Ataques «\emph{Attempted}»: intentos que no llegan a completarse y que
        aparecen etiquetados como ataque aunque no intercambien datos de
        aplicación.
  \item \emph{Shortcut learning}: atributos correlacionados con la etiqueta por
        accidente del montaje experimental, señaladamente las direcciones IP, que
        permiten al modelo acertar sin aprender nada del ataque.
\end{enumerate}

\subsection*{Dos de ellos son el objeto de esta memoria}

Los dos últimos no se citan aquí como contexto, sino que se miden en los
capítulos siguientes, y la conexión se puede enunciar antes de llegar a ellos.

El \emph{shortcut learning} se cuantifica en la sección~\ref{sec:fuga}. Añadir
las direcciones y el puerto de origen a la vista de metadatos produce un macro-F1
de $1.0000$ entre días distintos en cinco de los seis pares posibles, porque 15
de las 18 direcciones atacantes del conjunto empiezan por \texttt{18.} y las 12
víctimas están en \texttt{172.31.69.x}. Y se añade un resultado que Engelen no
reporta: la validación entre días no detecta el atajo, lo premia. Solo el cambio
de conjunto de datos lo rompe.

Las etiquetas «\emph{Attempted}» corresponden a lo que la
sección~\ref{sec:limite} mide como ceguera del análisis de contenido. En el día
de la Infiltration, el 94.4\,\% de las conexiones del reconocimiento no lleva un
solo byte de aplicación, y el ataque de fuerza bruta contra FTP del 14 de febrero
deja cero flujos analizables. No es que el enfoque funcione peor sobre esos
ataques: es que no tiene nada que leer.

\subsection*{Un defecto añadido por este trabajo}

A los cuatro anteriores, el capítulo~\ref{ch:regimenes} añade un quinto que la
literatura consultada no recoge: la evaluación dentro de un mismo día está
saturada. Medidas tres representaciones distintas con la misma métrica y el mismo
protocolo sobre seis días, ninguna baja de $0.9186$ y el análisis de contenido
no baja de $0.9734$. La comparación no discrimina, de modo que un resultado
alto obtenido así no informa sobre la capacidad de detección.

Esto tiene una consecuencia sobre cómo leer la literatura, y es que buena parte
de las cifras publicadas sobre estos conjuntos de datos se obtienen exactamente
con ese protocolo. No implica que esos trabajos estén equivocados, sino que sus
diferencias de centésimas no son necesariamente diferencias reales.

\section{Posicionamiento de este trabajo}
\label{sec:sota-posicionamiento}

De la literatura se toman las arquitecturas y las herramientas. La
representación de la carga útil como distribución de bytes procede de la
tradición de PAYL y Anagram \cite{Wang2004, Wang2006}; el byte-CNN sigue la línea
de \emph{Deep Packet} \cite{Lotfollahi2020}; el \emph{transformer} y la atención
cruzada se implementan siguiendo sus referencias respectivas; el escalado de
Platt \cite{Platt1999} y el error de calibración esperado \cite{Guo2017} son
estándar; y la crítica a los conjuntos de datos parte de Engelen \emph{et al.}
\cite{Engelen2021}.

Lo que este trabajo aporta no es una arquitectura, sino un conjunto de mediciones
sobre las condiciones en que esas arquitecturas se evalúan:

\begin{itemize}
  \item Una \emph{ground-truth} verificada sobre las capturas para siete días,
        que corrige la documentación oficial en tres de ellos
        (sección~\ref{sec:ground-truth}).
  \item La constatación de que la evaluación intra-día está saturada y no
        discrimina entre representaciones, con la hipótesis de la identidad del
        host contrastada y refutada (secciones~\ref{sec:saturacion}
        y~\ref{sec:por-que-satura}).
  \item La medición de que el éxito del análisis de contenido suele ser una
        huella de la herramienta y no del ataque, demostrada por
        interpretabilidad y por evasión (secciones~\ref{sec:saliency}
        y~\ref{sec:evasion}).
  \item El diagnóstico de que, entre dominios, lo que se rompe primero es la
        calibración y no la representación, con una receta que cuesta 25
        etiquetas (sección~\ref{sec:cross-dia}).
  \item La constatación de que la validación entre días no detecta la fuga de
        identificadores (sección~\ref{sec:fuga}).
\end{itemize}

A ello se añade una conclusión que el capítulo~\ref{ch:generalizacion} sostiene
con datos: sobre estos conjuntos, ninguna de las arquitecturas avanzadas
contrastadas mejora los resultados, y dos de ellas cuestan entre 2 y 49 veces
más. El cuello de botella no está en la capacidad del modelo, sino en la calidad
de la etiqueta y en el protocolo de medición. La
sección~\ref{sec:aportaciones} desarrolla este posicionamiento.
